\section*{Chapter \ref{ch:b1:stabilisation-policy} --- Stabilisation Policy and Its Limits}
% Numbers recomputed in code/bachelor-1/26-stabilisation-policy/exercises/solutions.py (tested).

\begin{solution}{exo:b1:stabilisation-policy:three}
Recognition, January to April; decision, April to September;
implementation, September to January (together the \omterm{def:b1:stabilisation-policy:lags}{inside lag}, a year);
the \omterm{def:b1:stabilisation-policy:lags}{outside lag}, January to the summer.
\end{solution}

\begin{solution}{exo:b1:stabilisation-policy:read}
Per point of growth lost, Germany's balance fell by $3.4/8.5 = 0.40$ point and
Ireland's by $14.2/10.4 = 1.37$. Stabilisers of similar size cannot explain
the gap: Ireland's tax receipts depended heavily on property and
construction, which collapsed, and the state supported its banks.
\end{solution}

\begin{solution}{exo:b1:stabilisation-policy:taylor}
$i = 4 + 0.5(-1) + 0.5(4 - 2) + 2 = 6.5\%$.
\end{solution}

\begin{solution}{exo:b1:stabilisation-policy:sign}
Negative: tax receipts fall with incomes and spending on benefits rises,
so the balance falls without any decision.
\end{solution}

\begin{solution}{exo:b1:stabilisation-policy:cycle-ex}
$1 + 0.25 - \cos(2\pi k/8)$: for $k = 1$, $1.25 - 0.707 = 0.543$; for $k = 2$, $1.25 - 0 = 1.25$; for $k = 4$,
$1.25 + 1 = 2.25$. Only the shortest lag helps.
\end{solution}

\begin{solution}{exo:b1:stabilisation-policy:uncertain}
$p^* = 4/1.5 = 2.667$. Expected losses: $(-4 + 2.667)^2 + 0.5 \times 2.667^2 = 5.333$; full response
$p = 4$: $0 + 0.5 \times 16 = 8$; no response: $16$.
\end{solution}

\begin{solution}{exo:b1:stabilisation-policy:slope}
The slope would be fixed, equal to the response of the budget balance per
point of output, the same everywhere, and the points would lie close to
the line. The measured slope mixes stabilisers with decisions (bank
rescues, stimulus), with tax systems of different size, and with
falls in revenue tied to particular sectors; and a change in growth is not
a change in the output gap. A weak slope does not mean weak stabilisers.
\end{solution}

\begin{solution}{exo:b1:stabilisation-policy:realtime}
The gap term is $0.5 \times (-3)$ instead of $0.5 \times (-1)$: the rate is $1$ point too low. Policy
is too loose, unemployment is pushed below its natural rate, and with
\omterm{def:b1:aggregate-supply:adaptive}{adaptive expectations} inflation rises year after year
(\cref{prop:b1:aggregate-supply:accelerate}).
\end{solution}

\begin{solution}{exo:b1:stabilisation-policy:local}
Smaller: the national central bank raises its rate in response to national
spending, and national taxes must pay for it; neither happens for one
region's spending relative to another's. Larger: a region spends much of
its income on goods from other regions, a leakage that disappears at the
national level.
\end{solution}

\begin{solution}{exo:b1:stabilisation-policy:ar}
$y_t = (\rho - \theta)y_{t-1} + \varepsilon_t$, so with $|\rho - \theta| < 1$, $\operatorname{var}(y) = \sigma^2/[1 - (\rho - \theta)^2]$. Without
policy it is $\sigma^2/(1 - \rho^2)$. Policy raises the variance when $|\rho - \theta| > \rho$, that is,
for $\theta > 0$, when $\theta > 2\rho$: a response too strong for the persistence of the gap
overshoots.
\end{solution}

\begin{solution}{exo:b1:stabilisation-policy:rule-disc}
The rule's expected loss is $9q$, \omterm{def:b1:stabilisation-policy:rule}{discretion}'s $4$: the rule is better when
$q < 4/9 = 0.44$, for example at $q = 0.2$ ($1.8$ against $4$), and worse at $q = 0.5$ ($4.5$).
\end{solution}

\begin{solution}{exo:b1:stabilisation-policy:refute}
Estimates from periods at the \omterm{def:b1:is-lm:liquidity-trap}{zero lower bound}, with spending changes
not caused by the state of the economy, giving multipliers no larger than
in normal times; for example narrative shocks near zero rates with
multipliers below 1. Ramey and Zubairy's mixed results at the bound are
the kind of evidence that bears on it.
\end{solution}

\begin{solution}{pb:b1:stabilisation-policy:brenland}
\textbf{1.} $-5$, $-3$, $-1.8$, $-1.08$, $-0.648$, $-0.389$, $-0.233$, $-0.14$.

\textbf{2.} $39.051$.

\textbf{3.} $-2$, $-1.2$, $-0.72$, $-0.432$, $-0.259$, $-0.156$, $-0.093$, $-0.056$.

\textbf{4.} $3$ points.

\textbf{5.} In April of the year before, in 5 of 153 recessions.

\textbf{6.} $-5$, $-3$, $0.7$, $0.42$, $0.252$, $0.151$, $0.091$, $0.054$.

\textbf{7.} $34.764$, against $39.051$: a gain, but policy overshoots from quarter 3.

\textbf{8.} $54.324$ with policy, $107.197$ without.

\textbf{9.} $31.25$ with policy, $25$ without.

\textbf{10.} The gap has closed by itself when the policy arrives: it pushes
output above potential by $2.5$ points.

\textbf{11.} How persistent the shock is, and how long the lag is.

\textbf{12.} $(30{,}000 - 12{,}000) \times 0.2 = 3{,}600$.

\textbf{13.} At 27,000, $3{,}000$; \omterm{def:b1:national-accounts-in-depth:gni}{disposable income} falls by $3{,}000 - 600 = 2{,}400$.

\textbf{14.} $20\%$: taxes absorb $600/3{,}000 = 20\%$ of the loss.

\textbf{15.} It becomes eligible for the transfer of at least 6,000 dalens,
which can more than offset the fall.

\textbf{16.} $1 - 1{,}532.3/2{,}226.9 = 0.312$: $31.2\%$.

\textbf{17.} Taxes $24.7\%$, transfers $6.5\%$.

\textbf{18.} $0.8 \times 1{,}532.3 = 1{,}225.8$ dalens per household, against $0.8 \times 2{,}226.9 = 1{,}781.5$.

\textbf{19.} It is built into the schedule: it responds to income as it is
earned.

\textbf{20.} $-5$, $-0.5$, $-0.3$, $-0.18$, $-0.108$, $-0.065$, $-0.039$, $-0.023$: loss $25.39$.

\textbf{21.} $41.237$.

\textbf{22.} $31.25$.

\textbf{23.} Lower than with a lag of two quarters for $\rho = 0.6$ ($25.39$ against $34.764$) and $\rho = 0.9$
($41.237$ against $54.324$); the same for $\rho = 0$.

\textbf{24.} It needs no decision: the \omterm{def:b1:stabilisation-policy:lags}{inside lag} is only the time to
measure the gap.

\textbf{25.} When shocks do not persist, any response to past gaps arrives
after the need for it and adds to the loss.

\textbf{26.} The income tax and the transfers, as incomes change.

\textbf{27.} Brenland's automatic stabilisers offset $31\%$ of a fall in market
income ($31.2\%$ of a $10\%$ fall in earnings): they need no forecast and no
decision, the two delays that make \omterm{def:b1:stabilisation-policy:rule}{discretion} risky.
\end{solution}
